\chapter{Interior positivity and deletion}

Throughout this chapter, the start \(j\) of a deleted block is zero-based:
the removed mathematical ranks are \(j+1,j+2,j+3\).  The hypotheses
\(1\leq j\) and \(j+3<n\) say exactly that all three removed ranks are
internal.

\begin{definition}\leanok
[Index embedding after deletion]
  \label{def:delete_index}
  \lean{Ninegates.deleteIndex}
  % SOURCE: half-formal.qmd, equation (8) preceding #lem-delete (read from references/half-formal.qmd)
  % SOURCE QUOTE: “定义保留点数的嵌入 $\iota_j:\{1,\ldots,n-3\}\to\{1,\ldots,n\}$ 为 $$ \iota_j(r)= \begin{cases} r,&r<j,\\ r+3,&r\ge j. \end{cases} \tag{8} $$ 它的像恰好是 $\{1,\ldots,n\}\setminus\{j,j+1,j+2\}$。”
  \textit{Source: half-formal.qmd, equation (8) preceding Lemma
  \texttt{\#lem-delete}.}
  For a zero-based block start \(j\) satisfying \(j+3\leq n\), embed the
  retained indices of a vector of length \(n-3\) into the original vector
  by
  \[
    \iota_j(r)=
    \begin{cases}
      r,&r<j,\\
      r+3,&r\geq j.
    \end{cases}
  \]
\end{definition}

\begin{definition}\leanok
[Deleting three coordinates]
  \label{def:delete_three}
  \lean{Ninegates.deleteThree}
  \uses{def:counts,def:delete_index}
  % SOURCE: half-formal.qmd, text following equation (8) (read from references/half-formal.qmd)
  % SOURCE QUOTE: “对长度为 $n$ 的向量 $v$，定义删除算子 $\operatorname{del}_j(v)_r=v_{\iota_j(r)}$，其中 $1\le r\le n-3$。”
  \textit{Source: half-formal.qmd, text following equation (8).}
  For a count vector \(v\) of length \(n\), define
  \(\operatorname{del}_j(v)\) of length \(n-3\) by
  \[
    \operatorname{del}_j(v)_r=v_{\iota_j(r)}.
  \]
\end{definition}

\begin{lemma}

\leanok
[A zero total has zero terms]
  \label{lem:red_eq_zero_of_sum_eq_zero}
  \lean{Ninegates.eq_zero_of_sum_eq_zero}
  In a finite sum of natural numbers, if the total is \(0\), then each
  summand is \(0\).
\end{lemma}
\begin{proof}

\leanok
  Each summand is bounded above by the whole finite sum.  If the whole
  sum is \(0\), antisymmetry in \(\mathbb N\) forces the chosen summand
  to be \(0\).
\end{proof}

\begin{lemma}

\leanok
[Canceling a drawn tile from residue two]
  \label{lem:modEq_of_add_one_modEq_two}
  \lean{Ninegates.modEq_of_add_one_modEq_two}
  If \(a+1\equiv2\pmod3\), then \(a\equiv1\pmod3\).
\end{lemma}
\begin{proof}

\leanok
  Subtract the same residue \(1\) from both sides of the congruence.
\end{proof}

\begin{lemma}

\leanok
[Canceling a drawn tile from residue zero]
  \label{lem:modEq_of_add_one_modEq_zero}
  \lean{Ninegates.modEq_of_add_one_modEq_zero}
  If \(a+1\equiv0\pmod3\), then \(a\equiv2\pmod3\).
\end{lemma}
\begin{proof}

\leanok
  Rewrite \(0\) as \(2+1\) modulo \(3\), then cancel the common residue
  \(1\).
\end{proof}

\begin{lemma}

\leanok
[Unweighted contribution of one run]
  \label{lem:red_run_count_sum}
  \lean{Ninegates.run_count_sum}
  \uses{lem:sum_if_fin_val_eq}
  A run beginning at zero-based start \(j\) contributes exactly three
  tiles to the unweighted total, counted with its multiplicity.
\end{lemma}
\begin{proof}

\leanok
  Split the finite sum into the three coordinates \(j,j+1,j+2\).  The
  singleton finite-sum lemma evaluates those three terms and all other
  coordinates contribute \(0\).
\end{proof}

\begin{definition}

\leanok
[Total left of a cut]
  \label{def:left_total}
  \lean{Ninegates.leftTotal}
  \uses{def:counts}
  For a cut coordinate \(k\), \(A_k(c)\) is the total multiplicity of
  all ranks strictly left of \(k\).
\end{definition}

\begin{definition}

\leanok
[Total right of a cut]
  \label{def:right_total}
  \lean{Ninegates.rightTotal}
  \uses{def:counts}
  For a cut coordinate \(k\), \(B_k(c)\) is the total multiplicity of
  all ranks strictly right of \(k\).
\end{definition}

\begin{definition}

\leanok
[Weighted rank sum left of a cut]
  \label{def:left_rank_sum}
  \lean{Ninegates.leftRankSum}
  \uses{def:counts}
  For a cut coordinate \(k\), \(S^L_k(c)\) is the one-based weighted
  rank sum of the coordinates strictly left of \(k\).
\end{definition}

\begin{definition}

\leanok
[Weighted rank sum right of a cut]
  \label{def:right_rank_sum}
  \lean{Ninegates.rightRankSum}
  \uses{def:counts}
  For a cut coordinate \(k\), \(S^R_k(c)\) is the one-based weighted
  rank sum of the coordinates strictly right of \(k\).
\end{definition}

\begin{lemma}

\leanok
[Filtered totals are additive]
  \label{lem:filtered_total_add}
  \lean{Ninegates.filteredTotal_add}
  For any predicate \(P\) on the finite coordinate set, the sum of
  \(a_i+b_i\) over the coordinates satisfying \(P\) is the sum over
  \(a_i\) plus the sum over \(b_i\).
\end{lemma}
\begin{proof}

\leanok
  Distribute addition inside each selected summand, then split the finite
  sum into the two component sums.
\end{proof}

\begin{lemma}

\leanok
[Filtered rank sums are additive]
  \label{lem:filtered_rank_add}
  \lean{Ninegates.filteredRank_add}
  For any predicate \(P\), the weighted filtered rank sum of
  \(a_i+b_i\) is the sum of the two weighted filtered rank sums.
\end{lemma}
\begin{proof}

\leanok
  On selected coordinates distribute the weight over \(a_i+b_i\), and
  then split the finite sum.
\end{proof}

\begin{lemma}

\leanok
[Filtered total of a basis vector]
  \label{lem:filtered_total_basis}
  \lean{Ninegates.filteredTotal_basis}
  \uses{def:basis}
  A multiple of the basis vector at \(x\) contributes \(m\) to a filtered
  total exactly when \(x\) satisfies the filter, and contributes \(0\)
  otherwise.
\end{lemma}
\begin{proof}

\leanok
  If \(x\) satisfies the filter, the finite sum has the single nonzero
  coordinate \(x\).  If it does not, every selected coordinate has zero
  basis value.
\end{proof}

\begin{lemma}

\leanok
[Filtered rank sum of a basis vector]
  \label{lem:filtered_rank_basis}
  \lean{Ninegates.filteredRank_basis}
  \uses{def:basis}
  A multiple of the basis vector at \(x\) contributes the weighted value
  \((x+1)m\) to a filtered rank sum exactly when \(x\) satisfies the
  filter, and contributes \(0\) otherwise.
\end{lemma}
\begin{proof}

\leanok
  The only possible nonzero coordinate is \(x\).  When it is selected the
  weight is its one-based rank; otherwise no selected coordinate
  contributes.
\end{proof}

\begin{lemma}

\leanok
[Filtered total after a draw]
  \label{lem:filtered_total_add_basis}
  \lean{Ninegates.filteredTotal_add_basis}
  \uses{lem:filtered_total_add,lem:filtered_total_basis}
  Adding one tile at \(x\) increases a filtered total by \(1\) if \(x\)
  satisfies the filter, and by \(0\) otherwise.
\end{lemma}
\begin{proof}

\leanok
  Apply filtered additivity to split off the basis vector, then evaluate
  that filtered basis contribution.
\end{proof}

\begin{lemma}

\leanok
[Filtered rank sum after a draw]
  \label{lem:filtered_rank_add_basis}
  \lean{Ninegates.filteredRank_add_basis}
  \uses{lem:filtered_rank_add,lem:filtered_rank_basis}
  Adding one tile at \(x\) increases a filtered rank sum by the
  one-based rank of \(x\) if \(x\) satisfies the filter, and by \(0\)
  otherwise.
\end{lemma}
\begin{proof}

\leanok
  Split the draw as a basis vector and evaluate the basis contribution in
  the weighted filtered sum.
\end{proof}

\begin{lemma}

\leanok
[Filtered total after removing a pair]
  \label{lem:filtered_total_remove_pair_add_pair}
  \lean{Ninegates.filteredTotal_removePair_add_pair}
  \uses{def:remove_pair,def:basis,lem:filtered_total_add,lem:filtered_total_basis}
  If \(w_p\ge2\), then the filtered total of \(w-2e_p\), plus the two
  removed tiles when \(p\) is selected by the filter, equals the filtered
  total of \(w\).
\end{lemma}
\begin{proof}

\leanok
  Coordinatewise, adding two basis tiles at \(p\) back to \(w-2e_p\)
  recovers \(w\).  Apply filtered additivity and the filtered basis
  evaluation.
\end{proof}

\begin{lemma}

\leanok
[Filtered rank sum after removing a pair]
  \label{lem:filtered_rank_remove_pair_add_pair}
  \lean{Ninegates.filteredRank_removePair_add_pair}
  \uses{def:remove_pair,def:basis,lem:filtered_rank_add,lem:filtered_rank_basis}
  If \(w_p\ge2\), then the filtered rank sum of \(w-2e_p\), plus twice
  the weighted rank of \(p\) when \(p\) is selected, equals the filtered
  rank sum of \(w\).
\end{lemma}
\begin{proof}

\leanok
  Add back the two basis tiles at \(p\) coordinatewise.  The weighted
  filtered additivity and basis lemmas evaluate the correction term.
\end{proof}

\begin{lemma}

\leanok
[Total split at a cut]
  \label{lem:total_split_at}
  \lean{Ninegates.total_split_at}
  \uses{def:total,def:left_total,def:right_total,lem:sum_if_fin_val_eq}
  The total count is the sum of the counts strictly left of \(k\), at
  \(k\), and strictly right of \(k\).
\end{lemma}
\begin{proof}

\leanok
  Partition every coordinate into the trichotomy \(i<k\), \(i=k\), or
  \(k<i\).  The middle singleton sum evaluates to \(c_k\).
\end{proof}

\begin{lemma}

\leanok
[Rank sum split at a cut]
  \label{lem:rank_sum_split_at}
  \lean{Ninegates.rankSum_split_at}
  \uses{def:rank_sum,def:left_rank_sum,def:right_rank_sum,lem:sum_if_fin_val_eq}
  The rank sum is the weighted sum strictly left of \(k\), plus the
  contribution of \(k\), plus the weighted sum strictly right of \(k\).
\end{lemma}
\begin{proof}

\leanok
  Use the same left/equal/right partition as for totals, now with each
  coordinate multiplied by its one-based rank.  The middle singleton sum
  is \((k+1)c_k\).
\end{proof}

\begin{lemma}

\leanok
[Meldable total is divisible by three]
  \label{lem:meldable_total_mod}
  \lean{Ninegates.meldable_total_mod}
  \uses{def:meldable,def:total,lem:red_run_count_sum,lem:sum_modEq_zero}
  The total multiplicity of any meldable vector is congruent to
  \(0\pmod3\).
\end{lemma}
\begin{proof}

\leanok
  In a meldable decomposition, triplets contribute multiples of \(3\).
  Each run contributes three tiles by
  Lemma~\ref{lem:red_run_count_sum}.  Summing these divisible
  contributions gives total residue \(0\).
\end{proof}

\begin{lemma}

\leanok
[A run crossing a zero coordinate has zero count]
  \label{lem:run_eq_zero_of_zero_coord}
  \lean{Ninegates.run_eq_zero_of_zero_coord}
  \uses{def:meldable,lem:red_eq_zero_of_sum_eq_zero}
  In a meldable coordinate equation, if coordinate \(k\) has value \(0\),
  then every run whose support contains \(k\) has multiplicity \(0\).
\end{lemma}
\begin{proof}

\leanok
  At coordinate \(k\), the triplet contribution and all covering run
  contributions are natural numbers summing to \(0\).  Hence each
  covering run term is \(0\).
\end{proof}

\begin{lemma}

\leanok
[Left total contribution of one run across a zero cut]
  \label{lem:run_left_count_mod}
  \lean{Ninegates.run_left_count_mod}
  \uses{lem:red_run_count_sum}
  If every run crossing a cut \(k\) has multiplicity \(0\), then the
  contribution of any single run to the total strictly left of \(k\) is
  congruent to \(0\pmod3\).
\end{lemma}
\begin{proof}

\leanok
  A run wholly left of the cut contributes all three of its tiles; a run
  crossing the cut contributes \(0\) by hypothesis; a run wholly right of
  the cut contributes nothing to the left side.  Each case has residue
  \(0\).
\end{proof}

\begin{lemma}

\leanok
[Left rank contribution of one run across a zero cut]
  \label{lem:run_left_rank_mod}
  \lean{Ninegates.run_left_rank_mod}
  \uses{lem:run_weight_sum}
  If every run crossing a cut \(k\) has multiplicity \(0\), then the
  contribution of any single run to the rank sum strictly left of \(k\)
  is congruent to \(0\pmod3\).
\end{lemma}
\begin{proof}

\leanok
  A run wholly left of the cut contributes the weighted sum of three
  consecutive ranks, divisible by \(3\).  A crossing run has zero
  multiplicity, and a wholly right run contributes nothing.
\end{proof}

\begin{lemma}

\leanok
[Left total of a meldable vector with a zero cut]
  \label{lem:meldable_left_total_mod_zero}
  \lean{Ninegates.meldable_leftTotal_mod_of_zero}
  \uses{def:left_total,lem:run_eq_zero_of_zero_coord,lem:run_left_count_mod,lem:sum_modEq_zero}
  If a meldable vector has value \(0\) at \(k\), then its total strictly
  left of \(k\) is congruent to \(0\pmod3\).
\end{lemma}
\begin{proof}

\leanok
  In a meldable decomposition, all runs crossing \(k\) vanish.  Triplets
  on the left are multiples of \(3\), and each remaining run contribution
  to the left total is congruent to \(0\) by
  Lemma~\ref{lem:run_left_count_mod}.  Sum the congruences.
\end{proof}

\begin{lemma}

\leanok
[Left rank sum of a meldable vector with a zero cut]
  \label{lem:meldable_left_rank_mod_zero}
  \lean{Ninegates.meldable_leftRank_mod_of_zero}
  \uses{def:left_rank_sum,lem:run_eq_zero_of_zero_coord,lem:run_left_rank_mod,lem:sum_modEq_zero}
  If a meldable vector has value \(0\) at \(k\), then its weighted rank
  sum strictly left of \(k\) is congruent to \(0\pmod3\).
\end{lemma}
\begin{proof}

\leanok
  Crossing runs vanish.  Triplet contributions on the left are weighted
  multiples of \(3\), and each noncrossing run contributes rank sum
  divisible by \(3\).  The filtered finite sum therefore has residue
  \(0\).
\end{proof}

\begin{lemma}

\leanok
[Right total of a meldable vector with a zero cut]
  \label{lem:meldable_right_total_mod_zero}
  \lean{Ninegates.meldable_rightTotal_mod_of_zero}
  \uses{def:right_total,lem:meldable_total_mod,lem:meldable_left_total_mod_zero,lem:total_split_at}
  If a meldable vector has value \(0\) at \(k\), then its total strictly
  right of \(k\) is congruent to \(0\pmod3\).
\end{lemma}
\begin{proof}

\leanok
  Split the whole total into left, middle, and right parts.  The whole
  total and left total are both \(0\) modulo \(3\), and the middle value
  is \(0\), so the right total is also \(0\) modulo \(3\).
\end{proof}

\begin{lemma}

\leanok
[Right rank sum of a meldable vector with a zero cut]
  \label{lem:meldable_right_rank_mod_zero}
  \lean{Ninegates.meldable_rightRank_mod_of_zero}
  \uses{def:right_rank_sum,lem:meldable_rank_sum,lem:meldable_left_rank_mod_zero,lem:rank_sum_split_at}
  If a meldable vector has value \(0\) at \(k\), then its weighted rank
  sum strictly right of \(k\) is congruent to \(0\pmod3\).
\end{lemma}
\begin{proof}

\leanok
  Split the whole rank sum into left, middle, and right parts.  The whole
  rank sum and left rank sum are \(0\) modulo \(3\), and the middle
  contribution is \(0\), so the right rank sum has residue \(0\).
\end{proof}

\begin{theorem}\leanok
[Interior coordinates are positive]
  \label{thm:interior_positive}
  \lean{Ninegates.interior_positive}
  \uses{def:live_wait,lem:meldable_left_total_mod_zero,lem:meldable_right_total_mod_zero,lem:meldable_left_rank_mod_zero,lem:meldable_right_rank_mod_zero,lem:filtered_total_add_basis,lem:filtered_total_remove_pair_add_pair,lem:filtered_rank_add_basis,lem:filtered_rank_remove_pair_add_pair,lem:modEq_of_add_one_modEq_two,lem:modEq_of_add_one_modEq_zero}
  % SOURCE: half-formal.qmd, #lem-interior-positive (read from references/half-formal.qmd)
  % SOURCE QUOTE: “设 $n\ge6$ 且 $3\mid n$。若 $c$ 是非空听的 $n$ 面听清一色牌型，则对每个内部点数 $k\in\{2,\ldots,n-1\}$ 都有 $c_k\ge1$。”
  \textit{Source: half-formal.qmd, Lemma
  \texttt{\#lem-interior-positive}.}
  Let \(n\geq6\) with \(3\mid n\).  If \(c\) is a live all-rank wait,
  then \(c_k\geq1\) for every internal rank \(2\leq k\leq n-1\).
\end{theorem}
% SOURCE QUOTE PROOF: “固定一个内部点数 $k$，反设 $c_k=0$。把牌型分成点数小于 $k$ 的左部与点数大于 $k$ 的右部，并记两侧总张数为 $A=\sum_{i=1}^{k-1}c_i$ 与 $B=\sum_{i=k+1}^{n}c_i$。因为 $k$ 是内部点数，左右两侧都非空：左侧至少含点数 $1$，右侧至少含点数 $n$。因此可以分别选取进张 $y<k$ 与 $z>k$。先考虑摸入左侧点数 $y$。完整牌型 $c+e_y$ 在点数 $k$ 处的重数仍为 $0$。任何跨越左右边界的顺子都必须含有点数 $k$： - 从 $k-2$ 开始的顺子为 $(k-2,k-1,k)$； - 从 $k-1$ 开始的顺子为 $(k-1,k,k+1)$； - 从 $k$ 开始的顺子为 $(k,k+1,k+2)$。由于没有点数 $k$，这些顺子都不可能出现。因此，和牌分解中的每个面子都完全位于左侧或完全位于右侧，雀头也完全位于其中一侧。全局只有一个雀头，所以恰有一侧由“一个雀头加若干面子”组成，另一侧完全由面子组成。任何纯面子牌型的总张数都是 $3$ 的倍数；任何“一个雀头加若干面子”的总张数模 $3$ 等于 $2$。摸入 $y$ 后，左右总张数分别为 $A+1$ 与 $B$，故 $$(A+1,B)\equiv(2,0)\quad\text{或}\quad(0,2)\pmod3.$$ 分别减去进张贡献，得到 $$(A,B)\equiv(1,0)\quad\text{或}\quad(2,2)\pmod3. \tag{3}$$ 再考虑摸入右侧点数 $z$。同样因为点数 $k$ 缺失，左右不能通过顺子相连。这时左右总张数分别为 $A$ 与 $B+1$，所以 $$(A,B+1)\equiv(2,0)\quad\text{或}\quad(0,2)\pmod3,$$ 从而 $$(A,B)\equiv(2,2)\quad\text{或}\quad(0,1)\pmod3. \tag{4}$$ 式 (3) 与式 (4) 必须同时成立。比较两个二元剩余类集合，它们唯一的公共元素是 $(2,2)$，所以 $$A\equiv2\pmod3,\qquad B\equiv2\pmod3. \tag{5}$$ 现在重新考虑任意左侧进张 $y<k$。由式 (5)，摸牌后 $A+1\equiv0\pmod3$，而 $B\equiv2\pmod3$。因此左侧不可能含雀头，右侧必须含唯一的雀头。于是，对每个 $y<k$，左侧向量 $L+e_y=(c_1,\ldots,c_{k-1})+e_y$ 都可面子分解。如果左侧至少有两个不同点数，即 $k-1\ge2$，则进张 $1$ 与进张 $2$ 都存在。令 $S_L=\sum_{i=1}^{k-1}i\,c_i$ 为左侧原牌型的点数和。由于 $L+e_1$ 与 $L+e_2$ 都可面子分解， @lem-rank-sum 分别给出 $$S_L+1\equiv0\pmod3,\qquad S_L+2\equiv0\pmod3.$$ 两式相减得到 $1\equiv0\pmod3$，矛盾。因此 $$k-1\le1. \tag{6}$$ 对右侧作完全对称的论证。由式 (5)，任意右侧进张 $z>k$ 后，右侧总张数 $B+1$ 被 $3$ 整除，所以雀头必须在左侧，而右侧向量 $R+e_z=(c_{k+1},\ldots,c_n)+e_z$ 可面子分解。如果右侧至少有两个不同点数，即 $n-k\ge2$，就可取 $z=k+1$ 与 $z=k+2$。记 $S_R=\sum_{i=k+1}^{n}i\,c_i$。由 @lem-rank-sum， $$S_R+(k+1)\equiv0\pmod3,\qquad S_R+(k+2)\equiv0\pmod3,$$ 再次相减得到矛盾。因此 $$n-k\le1. \tag{7}$$ 由式 (6) 与式 (7)，$n=(k-1)+1+(n-k)\le3$，这与 $n\ge6$ 矛盾。”
\begin{proof}

\leanok
  Fix an internal \(k\) and suppose \(c_k=0\).  Let \(A\) and \(B\)
  be the total multiplicities strictly left and strictly right of \(k\).
  Since a run crossing the cut would contain rank \(k\), every complete
  decomposition after a draw separates into left and right pieces, with
  the unique pair on exactly one side.

  Drawing on the left shows
  \[
    (A,B)\equiv(1,0)\ \text{or}\ (2,2)\pmod3,
  \]
  and drawing on the right shows
  \[
    (A,B)\equiv(2,2)\ \text{or}\ (0,1)\pmod3.
  \]
  Hence \(A\equiv B\equiv2\pmod3\).  It follows that after every
  left-side draw the left piece is purely meldable.  If the left side had
  two ranks, drawing its first and second ranks would produce two meldable
  vectors whose rank sums differ by \(1\), contradicting
  Lemma~\ref{lem:meldable_rank_sum}.  Thus the left side has at most one
  rank.  The symmetric argument gives the same conclusion on the right.
  Therefore \(n\leq3\), contradicting \(n\geq6\).
\end{proof}

\begin{lemma}

\leanok
[Option fold over appended lists]
  \label{lem:foldlM_append_option}
  \lean{Ninegates.foldlM_append_option}
  For a fold in the \(\mathrm{Option}\) monad, folding over an appended
  list is the same as folding over the prefix and, if that succeeds,
  continuing the fold over the suffix from the produced state.
\end{lemma}
\begin{proof}

\leanok
  Induct on the prefix.  The empty prefix is immediate.  For a nonempty
  prefix, expose the first monadic step; if it fails both sides fail, and
  if it succeeds the induction hypothesis applies to the tail.
\end{proof}

\begin{lemma}

\leanok
[Deleting a successful singleton scan block]
  \label{lem:scan_delete_three_ones_block}
  \lean{Ninegates.scan_delete_three_ones_block}
  \uses{lem:foldlM_append_option,def:scan_step,def:scan,lem:scan_three_ones}
  If a scan of \(\mathit{pref}\,[1,1,1]\,\mathit{suff}\) from state
  \(\sigma\) succeeds with final state \(\tau\), then the scan of
  \(\mathit{pref}\,\mathit{suff}\) from \(\sigma\) also succeeds with
  final state \(\tau\).
\end{lemma}
\begin{proof}

\leanok
  Split the successful scan at the prefix using
  Lemma~\ref{lem:foldlM_append_option}.  The prefix cannot have failed.
  Inspect the three scans of the entries \(1,1,1\); since the full scan
  succeeds, all three steps succeed.  They form the length-three scan of
  the constant-one vector, so Lemma~\ref{lem:scan_three_ones} says that
  the state after the block is the same as the state before it.  The
  suffix therefore starts from the same state with or without the block.
\end{proof}

\begin{lemma}

\leanok
[List form of a three-coordinate deletion]
  \label{lem:ofFn_deleteThree_eq_take_drop}
  \lean{Ninegates.ofFn_deleteThree_eq_take_drop}
  \uses{def:delete_three,def:delete_index}
  The list underlying \(\operatorname{del}_j(d)\) is obtained from the
  list underlying \(d\) by taking the first \(j\) entries and appending
  the entries after position \(j+2\).
\end{lemma}
\begin{proof}

\leanok
  Compare the two lists element by element.  An index before \(j\) is
  embedded unchanged by \(\iota_j\), while an index at or after \(j\) is
  embedded by adding \(3\).  Out-of-range lookups are out of range on
  both sides.
\end{proof}

\begin{lemma}

\leanok
[Splitting off three singleton entries]
  \label{lem:ofFn_eq_take_ones_drop}
  \lean{Ninegates.ofFn_eq_take_ones_drop}
  \uses{def:counts}
  If the three coordinates \(j,j+1,j+2\) of \(d\) are all \(1\), then
  the list underlying \(d\) is its prefix before \(j\), followed by
  \([1,1,1]\), followed by its suffix after \(j+2\).
\end{lemma}
\begin{proof}

\leanok
  Start from the identity \(l=\operatorname{take}_j(l)\) followed by
  \(\operatorname{drop}_j(l)\).  Repeatedly expose the first element of
  the remaining drop at positions \(j,j+1,j+2\), and use the singleton
  hypotheses to identify those exposed entries as \(1\).
\end{proof}

\begin{lemma}

\leanok
[Total as the sum of the underlying list]
  \label{lem:total_eq_list_sum}
  \lean{Ninegates.total_eq_list_sum}
  \uses{def:total,def:counts}
  The finite total of a count vector equals the sum of the entries of its
  underlying list.
\end{lemma}
\begin{proof}

\leanok
  Induct on the vector length.  In the successor case, split the finite
  sum into the head coordinate and the sum over successor coordinates,
  and identify the list tail with the list of the shortened successor
  vector.
\end{proof}

\begin{lemma}

\leanok
[The retained-index embedding is injective]
  \label{lem:deleteIndex_injective}
  \lean{Ninegates.deleteIndex_injective}
  \uses{def:delete_index}
  The map \(\iota_j\) from retained indices into the original coordinate
  set is injective.
\end{lemma}
\begin{proof}

\leanok
  Split two retained indices according to whether they lie before the
  deleted block.  On the left side \(\iota_j\) is the identity; on the
  right side it is translation by \(3\); and the mixed cases cannot have
  equal images because the deleted block separates the two ranges.
\end{proof}

\begin{lemma}

\leanok
[Retained indices avoid the deleted block]
  \label{lem:deleteIndex_not_deleted}
  \lean{Ninegates.deleteIndex_not_deleted}
  \uses{def:delete_index}
  No value in the image of \(\iota_j\) lies in the deleted interval
  \([j,j+3)\).
\end{lemma}
\begin{proof}

\leanok
  If the retained index is before \(j\), its image is before the deleted
  interval.  Otherwise its image is at least \(j+3\).
\end{proof}

\begin{lemma}

\leanok
[Every nondeleted index has a retained preimage]
  \label{lem:deleteIndex_preimage_of_not_deleted}
  \lean{Ninegates.deleteIndex_preimage_of_not_deleted}
  \uses{def:delete_index}
  Every original index outside the deleted interval \([j,j+3)\) is in
  the image of \(\iota_j\).
\end{lemma}
\begin{proof}

\leanok
  If the original index is before \(j\), use the same value as the
  retained preimage.  Otherwise, being outside the deleted interval means
  it is at least \(j+3\), so subtract \(3\) and use that retained index.
\end{proof}

\begin{lemma}

\leanok
[Basis vectors commute with retained-index embedding]
  \label{lem:basis_deleteIndex}
  \lean{Ninegates.basis_deleteIndex}
  \uses{def:basis,def:delete_index,lem:deleteIndex_injective}
  Pulling a basis vector back along \(\iota_j\) gives the corresponding
  basis vector on the retained coordinate set.
\end{lemma}
\begin{proof}

\leanok
  At the distinguished retained coordinate both sides are \(1\).  At any
  other retained coordinate, injectivity of \(\iota_j\) shows that the
  embedded coordinates are distinct, so both basis values are \(0\).
\end{proof}

\begin{lemma}

\leanok
[Deleting after a retained draw and retained pair removal]
  \label{lem:deleteThree_removePair_add_basis}
  \lean{Ninegates.deleteThree_removePair_add_basis}
  \uses{def:delete_three,def:remove_pair,lem:basis_deleteIndex}
  Deleting after drawing a retained tile \(x'\) and removing a retained
  pair \(p'\) agrees coordinatewise with first deleting and then drawing
  \(x'\) and removing \(p'\) in the shortened coordinate set.
\end{lemma}
\begin{proof}

\leanok
  Evaluate both vectors at an arbitrary retained coordinate.  The
  underlying hand is compared by the definition of deletion, pair removal
  is evaluated at the embedded pair coordinate, and the basis-vector
  contribution is transported by Lemma~\ref{lem:basis_deleteIndex}.
\end{proof}

\begin{lemma}\leanok
[Deleting three singletons preserves meldability]
  \label{lem:delete_three_meldable}
  \lean{Ninegates.delete_three_meldable}
  \uses{def:delete_three,thm:meldable_iff_scan,lem:ofFn_deleteThree_eq_take_drop,lem:ofFn_eq_take_ones_drop,lem:scan_delete_three_ones_block}
  % SOURCE: half-formal.qmd, proof of #lem-delete, scan-preservation paragraph (read from references/half-formal.qmd)
  % SOURCE QUOTE: “现在要证明 $\operatorname{del}_j(d)$ 可面子分解。由 $d$ 可面子分解以及 @lem-recurrence，它的规范扫描从 $(0,0)$ 出发，不会失败，并在读完全部 $n$ 个坐标后回到 $(0,0)$。由于进张与雀头都不在被删区间，$d$ 在三个被删坐标上的重数仍为 $d_j=d_{j+1}=d_{j+2}=1$。设扫描在读入坐标 $j$ 之前的状态为 $\sigma$。原扫描能够顺利读完上述三个 $1$，所以 @lem-three-ones-state 适用，并说明读完坐标 $j,j+1,j+2$ 后状态仍为 $\sigma$。删除这三个坐标后，缩短向量的扫描在读完原坐标 $j-1$ 时同样处于 $\sigma$；它接下来直接读取原坐标 $j+3$。这与原扫描读完三个 $1$ 后读取 $j+3$ 时的状态完全相同。因此，从该处直到末尾，缩短扫描与原扫描逐步产生相同的状态。原扫描不失败并最终回到 $(0,0)$，所以缩短扫描也不失败并最终回到 $(0,0)$。再次使用 @lem-recurrence，得到 $\operatorname{del}_j(d)$ 可面子分解。”
  \textit{Source: half-formal.qmd, proof of Lemma
  \texttt{\#lem-delete}.}
  If \(d\) is meldable and its three coordinates beginning at \(j\) are
  all \(1\), then \(\operatorname{del}_j(d)\) is meldable.
\end{lemma}
% SOURCE QUOTE PROOF: “设扫描在读入坐标 $j$ 之前的状态为 $\sigma$。原扫描能够顺利读完上述三个 $1$，所以 @lem-three-ones-state 适用，并说明读完坐标 $j,j+1,j+2$ 后状态仍为 $\sigma$。删除这三个坐标后，缩短向量的扫描在读完原坐标 $j-1$ 时同样处于 $\sigma$；它接下来直接读取原坐标 $j+3$。这与原扫描读完三个 $1$ 后读取 $j+3$ 时的状态完全相同。因此，从该处直到末尾，缩短扫描与原扫描逐步产生相同的状态。原扫描不失败并最终回到 $(0,0)$，所以缩短扫描也不失败并最终回到 $(0,0)$。”
\begin{proof}

\leanok
  By Theorem~\ref{thm:meldable_iff_scan}, the scan of \(d\) succeeds from
  zero to zero.  Let \(\sigma\) be the state just before the three
  singleton coordinates.  Lemma~\ref{lem:scan_three_ones} says that the
  state just after them is again \(\sigma\).  The shortened scan has the
  same prefix state and then reads exactly the same suffix, so it also
  succeeds and finishes at zero.  Apply
  Theorem~\ref{thm:meldable_iff_scan} again.
\end{proof}

\begin{theorem}\leanok
[Deletion preserves a live all-rank wait]
  \label{thm:delete_three_live_wait}
  \lean{Ninegates.delete_three_live_wait}
  \uses{def:live_wait,def:delete_three,lem:delete_three_meldable,lem:total_eq_list_sum,lem:ofFn_deleteThree_eq_take_drop,lem:ofFn_eq_take_ones_drop,lem:deleteIndex_not_deleted,lem:deleteIndex_preimage_of_not_deleted,lem:basis_deleteIndex,lem:deleteThree_removePair_add_basis}
  % SOURCE: half-formal.qmd, #lem-delete (read from references/half-formal.qmd)
  % SOURCE QUOTE: “设 $n\ge9$ 且 $3\mid n$，并设 $c$ 是非空听的 $n$ 面听牌型。假设存在 $j$ 满足 $2\le j\le n-3$ 以及 $c_j=c_{j+1}=c_{j+2}=1$。令 $c'=\operatorname{del}_j(c)$。则 $c'$ 是非空听的 $(n-3)$ 面听牌型。”
  \textit{Source: half-formal.qmd, Lemma \texttt{\#lem-delete}.}
  Let \(n\geq9\), \(3\mid n\), and let \(c\) be a live all-rank wait.
  If the three internal coordinates beginning at \(j\) all have
  multiplicity \(1\), then \(\operatorname{del}_j(c)\) is a live
  all-rank wait of length \(n-3\).
\end{theorem}
% SOURCE QUOTE PROOF: “先验证 $c'$ 满足牌数与非空听的基本条件。由三个被删坐标的重数均为 $1$，它们总共恰有三张牌，因此 $$\sum_{r=1}^{n-3}c'_r=\sum_{i=1}^{n}c_i-(c_j+c_{j+1}+c_{j+2})=(n+4)-3=(n-3)+4.$$ 又因为 $c'$ 的每个坐标都直接取自 $c$ 的一个保留坐标，而 $0\le c_i\le3$，所以对所有 $r$ 都有 $0\le c'_r\le3$。特别地，$c'$ 中每一种候选进张都不是空听。剩下只需证明：对每个新点数 $x'\in\{1,\ldots,n-3\}$，摸入 $x'$ 后都能和牌。令 $x=\iota_j(x')$。这里 $x$ 是原牌型中与 $x'$ 对应的保留点数，所以 $x\notin\{j,j+1,j+2\}$。因为 $c$ 是 $n$ 面听，$c+e_x$ 是完整牌型。故存在某个雀头点数 $p\in\{1,\ldots,n\}$，使得 $d=c+e_x-2e_p$ 可面子分解。我们先证明雀头不在被删区间内。由假设，被删区间中的每个点数原本都只有一张；进张 $x$ 又不在该区间，所以在 $c+e_x$ 中仍有 $(c+e_x)_j=(c+e_x)_{j+1}=(c+e_x)_{j+2}=1$。雀头需要同一点数的两张牌，因此 $p\notin\{j,j+1,j+2\}$。因此存在唯一的新点数 $p'\in\{1,\ldots,n-3\}$ 满足 $\iota_j(p')=p$。而且删除算子与这次进张、雀头删除相容： $$\operatorname{del}_j(d)=c'+e_{x'}-2e_{p'}. \tag{9}$$ 式 (9) 可以逐坐标验证：在 $x'$ 坐标增加一张，在 $p'$ 坐标删除两张，其余保留坐标不变。现在要证明 $\operatorname{del}_j(d)$ 可面子分解。由 $d$ 可面子分解以及 @lem-recurrence，它的规范扫描从 $(0,0)$ 出发，不会失败，并在读完全部 $n$ 个坐标后回到 $(0,0)$。由于进张与雀头都不在被删区间，$d$ 在三个被删坐标上的重数仍为 $d_j=d_{j+1}=d_{j+2}=1$。设扫描在读入坐标 $j$ 之前的状态为 $\sigma$。原扫描能够顺利读完上述三个 $1$，所以 @lem-three-ones-state 适用，并说明读完坐标 $j,j+1,j+2$ 后状态仍为 $\sigma$。删除这三个坐标后，缩短向量的扫描在读完原坐标 $j-1$ 时同样处于 $\sigma$；它接下来直接读取原坐标 $j+3$。这与原扫描读完三个 $1$ 后读取 $j+3$ 时的状态完全相同。因此，从该处直到末尾，缩短扫描与原扫描逐步产生相同的状态。原扫描不失败并最终回到 $(0,0)$，所以缩短扫描也不失败并最终回到 $(0,0)$。再次使用 @lem-recurrence，得到 $\operatorname{del}_j(d)$ 可面子分解。结合式 (9)，$c'+e_{x'}$ 可以分解为雀头 $2e_{p'}$ 与若干面子，所以它是完整牌型。由于 $x'$ 是任意的新点数，$c'$ 对所有 $n-3$ 种进张都能和牌，故 $c'$ 是非空听的 $(n-3)$ 面听。”
\begin{proof}

\leanok
  The deleted coordinates contain exactly three tiles, so the shortened
  total is \((n+4)-3=(n-3)+4\).  Every retained coordinate still has
  multiplicity at most \(3\).

  Fix a new draw \(x'\), and let \(x=\iota_j(x')\).  Choose a complete
  decomposition of \(c+e_x\), with pair at \(p\) and meldable remainder
  \(d\).  Since \(x\) is retained, the three deleted coordinates still
  have multiplicity \(1\), so the pair cannot lie among them.  Hence
  \(p=\iota_j(p')\) for a unique retained \(p'\).  Coordinatewise
  calculation gives
  \[
    \operatorname{del}_j(d)
      =\operatorname{del}_j(c)+e_{x'}-2e_{p'}.
  \]
  The left side is meldable by
  Lemma~\ref{lem:delete_three_meldable}.  Thus every draw completes the
  shortened hand, proving the live-wait condition.
\end{proof}

\begin{theorem}\leanok
[Deletion detects the standard hand]
  \label{thm:delete_three_standard}
  \lean{Ninegates.delete_three_standard_iff}
  \uses{def:standard,def:delete_three}
  % SOURCE: half-formal.qmd, final assertion of #lem-delete (read from references/half-formal.qmd)
  % SOURCE QUOTE: “此外，$c=(3,1,\ldots,1,3)$ 当且仅当 $c'=(3,1,\ldots,1,3)$。”
  \textit{Source: half-formal.qmd, final assertion of Lemma
  \texttt{\#lem-delete}.}
  If the deleted block is internal and consists of three singleton
  coordinates, then
  \[
    c=G_n\quad\Longleftrightarrow\quad
    \operatorname{del}_j(c)=G_{n-3}.
  \]
\end{theorem}
% SOURCE QUOTE PROOF: “因为 $j\ge2$ 且 $j+2\le n-1$，被删的三个坐标都是内部坐标，两个端点均被保留。删除的三个重数又全部等于 $1$。因此： - 如果 $c=(3,1,\ldots,1,3)$，删除三个内部的 $1$ 后显然仍有 $c'=(3,1,\ldots,1,3)$； - 如果 $c'=(3,1,\ldots,1,3)$，把三个重数为 $1$ 的内部坐标插回原位置后，两端仍为 $3$，全部内部坐标仍为 $1$，所以 $c=(3,1,\ldots,1,3)$。”
\begin{proof}

\leanok
  The deleted positions are internal, both endpoints are retained, and the
  removed multiplicities are all \(1\).  Deleting them from \(G_n\)
  therefore gives \(G_{n-3}\).  Conversely, if the retained vector is
  \(G_{n-3}\), reinserting three internal \(1\)'s leaves endpoint
  multiplicities \(3\) and makes every internal multiplicity \(1\), so the
  original vector is \(G_n\).
\end{proof}
